% !TEX root = ../macro.tex
\section{Endogenous Growth and Resource Allocation}\label{sec:endogenous-growth}

Intermediate macroeconomics introduces AK, education, and innovation as ways
to explain sustained growth. This chapter asks what closes those explanations:
who accumulates the relevant stock, what it costs, which return guides the
choice, and whether that return is the one society receives. We reuse Ramsey's
optimization tools rather than treating endogenous growth as an unrelated list
of production functions.

\subsection{From intermediate macroeconomics and Ramsey to a growth engine}

\subsubsection{What the familiar diagrams leave to be explained}

A machine, a trained worker, and a usable design can all raise output, but they
are not accumulated in the same way. Building equipment uses goods. Training
uses time and other resources. Developing a design can be costly even when
another firm can subsequently use the same instructions without taking them
away from the first. These differences affect both the resource constraint
and the incentive to invest.

It is tempting to put each item into a production function and declare growth
explained. Instead, ask which scarcity remains after accumulation. If machines
are added while skills or usable methods remain fixed, returns may fall. If
skills and equipment can grow together, the relevant return might not fall.
If ideas can be shared, the social gain from discovery may exceed what the
inventor can charge. We need to identify that mechanism before choosing the
law of motion or discussing policy.

The familiar intermediate-level statements are therefore starting points.
``Saving can affect growth'' needs a technology that sustains returns and a
rule determining saving. ``Education raises productivity'' needs a skill
accumulation technology and a time allocation. ``Research creates ideas''
needs a resource cost and a reason someone pays it. The models below fill
these particular gaps; they do not repeat the Solow derivation.

\subsubsection{Ramsey endogenizes saving, not automatically long-run growth}

With fixed technology and population, the Ramsey benchmark has
\begin{equation}
\dot k=f(k)-c-\delta k,\qquad
\frac{\dot c}{c}=\frac{f'(k)-\delta-\rho}{\theta}.
\end{equation}
Saving is chosen optimally. Yet with diminishing marginal product, accumulating
$k$ lowers its return. At the stationary optimum,
$f'(k^*)=\delta+\rho$ and consumption stops growing. In the earlier version
with exogenous labor-efficiency growth, consumption per person can grow, but
that long-run technology growth rate is still supplied from outside the model.

The next question is thus about the return to \emph{all reproducible inputs},
not whether households have an Euler equation. Keep the household problem,
change the accumulation opportunities, and check what replaces the stationary
capital target. OLG changed who lives, saves, and receives transfers. Here we
return explicitly to a continuing Ramsey household to isolate a different
change: the mechanism that makes productive capacity grow.

\Needspace{10\baselineskip}
\begin{blueblock}
\textbf{Endogenous saving and endogenous growth are different claims.}
An optimal saving decision can lead to a stationary economy. Conversely, an
AK equation with an imposed saving share can generate a growth rate without
explaining saving. To obtain both, combine a specified accumulation technology
with household optimization, market or planner closure, and the appropriate
terminal conditions. Changing notation from an exogenous $g$ to a function
called $g$ does not provide those mechanisms.
\end{blueblock}

\subsection{AK: retain Ramsey's choice, change the return}

\Needspace{12\baselineskip}
\subsubsection{Complete conditions for the one-stock benchmark}

\begin{modelsetup}[1]
\begin{modelcondition}{01}{Time, population, goods, and states}
Time is continuous. A closed economy has one good, no government in the
baseline, no uncertainty, and a fixed population normalized to one. A
continuing household owns the given stock $k(0)=k_0>0$ and firms.
Consequently aggregate and per-person quantities coincide in this benchmark.
There is no age structure, leisure choice, or exogenous technology growth.
\end{modelcondition}
\begin{modelcondition}{02}{Technology and accumulation}
Production is linear in a reproducible stock:
\[
y=Bk,\qquad B>0\text{ constant},\qquad
\dot k=(B-\delta)k-c,\quad \delta\geq0.
\]
The stock can be interpreted as a broad productive composite only if its
components can be reproduced with the assumed technology. Goods can be
converted into the stock at unit cost. Strictly diminishing returns and the
Solow Inada limits are deliberately replaced; this is not Cobb--Douglas
with the same fixed factor left implicit.
\end{modelcondition}
\begin{modelcondition}{03}{Preferences and household choice}
The household chooses a positive consumption path and feasible $k(t)\geq0$:
\[
\max\int_0^\infty e^{-\rho t}u(c(t))\,dt,\qquad
u'(c)=c^{-\theta},\quad\rho>0,\quad\theta>0.
\]
Use $u(c)=c^{1-\theta}/(1-\theta)$ for $\theta\ne1$ and $u(c)=\log c$ for
$\theta=1$. Saving is the residual of output after consumption, not an
independently imposed share.
\end{modelcondition}
\begin{modelcondition}{04}{Markets, ownership, and returns}
In the competitive interpretation, firms rent the productive stock at
$q=B$ and earn zero economic profit after paying it. There is no separate
productive labor input in this one-stock technology. The household owns the
stock, receives the net return $r=B-\delta$, and can adjust investment without
costs. Without externalities or other wedges, its budget coincides with the
aggregate resource constraint and implements the planner allocation.
\end{modelcondition}
\begin{modelcondition}{05}{Admissibility and the growing solution}
The planner's transversality condition is
\[
\lim_{t\to\infty}e^{-\rho t}c(t)^{-\theta}k(t)=0.
\]
For the interior positive-growth case studied below, impose
\[
g\equiv\frac{B-\delta-\rho}{\theta}>0,\qquad
\xi\equiv\rho-(1-\theta)g>0.
\]
The second restriction ensures finite allocation-dependent utility and the
terminal condition along the proposed balanced growth path. For $\theta<1$
it restricts how productive accumulation can be relative to discounting;
a formal Euler solution alone need not be a well-posed optimum.
\end{modelcondition}
\end{modelsetup}

\subsubsection{The same Hamiltonian now gives a constant growth rate}

The current-value Hamiltonian and its control and costate conditions are
\begin{align}
\mathcal H&=u(c)+\mu[(B-\delta)k-c],\\
\mu&=c^{-\theta},\qquad
\dot\mu=[\rho-(B-\delta)]\mu.
\end{align}
Since $\dot\mu/\mu=-\theta\dot c/c$,
\begin{equation}\label{eq:eg-ak-growth}
\frac{\dot c}{c}=\frac{B-\delta-\rho}{\theta}=g.
\end{equation}
The marginal product no longer falls as the stock rises. Waiting remains
rewarded at the same rate, so the optimal consumption growth rate need not
fade toward zero. This is the change in technology doing the work; the
intertemporal preference condition is still Ramsey's.

The Euler equation gives $c(t)=c_0e^{gt}$ but not $c_0$. Set
$a\equiv B-\delta$ and integrate $\dot k=ak-c_0e^{gt}$, using
$\xi=a-g=\rho-(1-\theta)g>0$:
\begin{align}
k(t)&=e^{at}\left[k_0-c_0\int_0^t e^{-(a-g)v}\,dv\right],\\
&=\left(k_0-\frac{c_0}{\xi}\right)e^{at}
  +\frac{c_0}{\xi}e^{gt}.
\end{align}
If $c_0>\xi k_0$, the first term eventually makes capital negative. If
$c_0<\xi k_0$, capital grows at the faster rate $a$ and its limiting discounted
utility value is positive: the household leaves valuable resources unused.
Feasibility and transversality therefore select
\begin{equation}\label{eq:eg-ak-path}
c_0=\xi k_0,\qquad k(t)=k_0e^{gt},\qquad c(t)=\xi k(t).
\end{equation}
Indeed, along this path the terminal expression is proportional to
$e^{-[\rho-(1-\theta)g]t}=e^{-\xi t}$ and vanishes.

\Needspace{9\baselineskip}
\begin{greenblock}
\textbf{A balanced growth path replaces the stationary capital target.}
There is no finite positive stationary $k$ on this growing solution.
Consumption jumps to the feasible optimal ratio $c/k=\xi$ for the inherited
stock, then both grow at $g$. The saving share is an outcome,
$s^*=(\delta+g)/B=1-\xi/B$, rather than Solow's input. Under the stated
positive-growth conditions it lies between zero and one.
\end{greenblock}

\subsubsection{A level change, a return change, and a policy wedge}

A one-time increase in initial capital, holding $B$ and preferences fixed,
raises the levels of capital and consumption but leaves $g$ unchanged.
A permanent change in $B$ changes the return and hence $g$; comparing the
new paths requires the new initial consumption ratio as well as the slope.
The absence of a transition in the ratio is a special feature of linear
technology and these preferences, not a general result for growth models.

For example, $B=0.10$, $\delta=0.02$, $\rho=0.04$, and $\theta=2$ imply
$g=0.02$, $c/k=0.06$, and $s^*=0.40$. Raising $k_0$ by 10\% raises the whole
path by 10\%. Raising $B$ to $0.12$ instead gives $g=0.03$ and $c/k=0.07$.
These are illustrative model numbers, not empirical return estimates.

A tax rate $\tau$ on net capital income, rebated as a lump-sum transfer,
leaves aggregate resources unchanged but changes the household's marginal
return. For an interior admissible growing equilibrium,
\begin{equation}
g_\tau=\frac{(1-\tau)(B-\delta)-\rho}{\theta},\qquad
\frac{c}{k}=B-\delta-g_\tau.
\end{equation}
The finite-welfare restriction $\rho-(1-\theta)g_\tau>0$ must still hold.
This is a permanent growth effect under the AK assumption. The corresponding
Ramsey tax experiment with diminishing returns changed stationary levels;
that difference comes from the accumulation technology. An empirical claim
that a tax changes trend growth would require evidence for the relevant
mechanism, not this formula alone.

\subsubsection{Private diminishing returns can coexist with aggregate AK}

The one-stock competitive benchmark need not contain a market failure. An
alternative construction does: suppose firm $i$ has
\begin{equation}
y_i=B k_i^\alpha(\bar k\ell_i)^{1-\alpha},\qquad 0<\alpha<1,
\end{equation}
where $\bar k=K/L$ is an economy-wide knowledge spillover associated with
capital, taken as given by each small firm. Labor satisfies $\sum_i\ell_i=1$.
Firms have constant returns in their private inputs $k_i,\ell_i$, and take
factor prices and $\bar k$ as given. In a symmetric allocation with the same
capital per worker across firms, aggregate output is $Y=BK$.

Differentiate the firm's own technology before imposing symmetry. Then
\begin{equation}
q_D=\alpha B,\qquad w=(1-\alpha)BK,\qquad
r_D=\alpha B-\delta.
\end{equation}
Factor payments still exhaust output. A planner who internalizes how $K$
raises the common spillover instead obtains marginal product $B$. With the
same household preferences and appropriate admissibility restrictions,
\begin{equation}
g_D=\frac{\alpha B-\delta-\rho}{\theta},\qquad
 g_P=\frac{B-\delta-\rho}{\theta}.
\end{equation}
Where both growing solutions are feasible, $g_D<g_P$. These are two different
interpretations of the same aggregate $Y=BK$: either $B$ is privately earned,
or part of the social return is external. Aggregate linearity alone cannot
tell us which policy problem is present. Learning during investment or shared
production experience can motivate a spillover, but its size has to be
specified and measured rather than inferred from the AK notation.

\subsection{Human capital: explain the reproducible composite}

\subsubsection{Why adding a second stock is not enough}

The augmented Solow model added physical and human capital with diminishing
returns to their joint accumulation. Its balanced growth still depended on
exogenous technology. To change that result, give skill formation its own
technology: time spent learning today affects the ability to produce and learn
tomorrow. The question is now how much time to divert from current output.
This motivates a two-sector Uzawa--Lucas-style benchmark. We omit average-skill
externalities to isolate the private accumulation mechanism.\footnote{Robert
E. Lucas Jr. (1988), ``On the Mechanics of Economic Development,''
\textit{Journal of Monetary Economics}, 22(1), 3--42, Section 4.
The simplified model here is not the full externality specification in that
paper. \url{https://doi.org/10.1016/0304-3932(88)90168-7}.}

\Needspace{12\baselineskip}
\begin{modelsetup}[1]
\begin{modelcondition}{01}{Household, stocks, and the scope of the experiment}
Keep the continuous-time, fixed-population household and CRRA preferences
above, with $K_0,H_0>0$. $H$ is a stock of skill in the continuing household,
not hours worked or the number of school years. The household chooses
consumption $C>0$ and a work fraction $v\in[0,1]$. The remaining fraction
$1-v$ is used for learning; leisure is omitted. There is no uncertainty,
knowledge spillover to other households, or public subsidy in this benchmark.
\end{modelcondition}
\begin{modelcondition}{02}{Production and the goods constraint}
Goods production is
\[
Y=K^\alpha(vH)^{1-\alpha},\qquad 0<\alpha<1,\qquad
\dot K=Y-C-\delta K.
\]
Goods can be consumed or accumulated as physical capital. Current production
uses only the skill-weighted time assigned to work. The technology is constant
returns in $K,H$ at fixed $v$, despite a diminishing marginal product of $K$
when $H$ is fixed. The household/planner controls the two uses of its resources.
\end{modelcondition}
\begin{modelcondition}{03}{Learning technology and its cost}
Skills accumulate according to
\[
\dot H=[\eta(1-v)-\delta_H]H,\qquad \eta>0,\quad\delta_H\geq0.
\]
Learning uses time, so its resource cost is forgone production, even though
there is no separate goods expenditure in this stripped-down specification.
A given learning fraction adds skill in proportion to the existing stock.
Skill depreciation is explicit; finite careers, cohort replacement, and
education quality differences are not separately modeled. Indefinite aggregate
skill accumulation is an assumption about this continuing household, not an
assertion that one person can study forever.
\end{modelcondition}
\begin{modelcondition}{04}{Optimality, boundaries, and balanced growth}
This is a planner/household production problem with no externality. The
Hamiltonian conditions below apply to an interior allocation $0<v<1$,
$K,H,C>0$, with terminal conditions
\[
\lim_{t\to\infty}e^{-\rho t}\lambda(t)K(t)=0,\qquad
\lim_{t\to\infty}e^{-\rho t}\mu(t)H(t)=0.
\]
For a positive-growth balanced path require
$g=(\eta-\delta_H-\rho)/\theta>0$ and
$\rho-(1-\theta)g>0$. They imply the interior work fraction derived below;
if the implied fraction lies outside $[0,1]$, the interior solution cannot
be used. Capital and skill are both inherited; an arbitrary initial ratio
need not already be on the balanced growth path.
\end{modelcondition}
\end{modelsetup}

\subsubsection{Time allocation and the two shadow values}

Write the current-value Hamiltonian as
\begin{equation}
\mathcal H=u(C)+\lambda[Y-C-\delta K]
 +\mu[\eta(1-v)-\delta_H]H.
\end{equation}
The consumption condition is $\lambda=C^{-\theta}$. The work-time condition
and the physical-capital costate equation are
\begin{align}
\lambda(1-\alpha)Y/v&=\mu\eta H,\label{eq:eg-human-time}\\
\dot\lambda/\lambda&=\rho+\delta-\alpha Y/K.
\end{align}
An extra unit of work time produces goods now; an extra unit of learning time
produces skills valued at $\mu$. The first equation compares those two uses
of time in the same utility units.

For the skill stock,
\begin{equation}
\dot\mu=\rho\mu-\lambda(1-\alpha)Y/H
 -\mu[\eta(1-v)-\delta_H].
\end{equation}
Use \eqref{eq:eg-human-time} to replace
$\lambda(1-\alpha)Y/H$ by $\mu\eta v$. Then
\begin{equation}\label{eq:eg-human-costate}
\frac{\dot\mu}{\mu}=\rho-\eta+\delta_H.
\end{equation}
The return to skill includes its contribution to further skill formation.
Looking only at its current wage contribution would miss that part of the
accumulation opportunity.

\subsubsection{Balanced growth and the work--learning split}

On a balanced path, $K,H,C$ grow at the same rate $g$ and $v$ is constant.
Equation~\eqref{eq:eg-human-time} then makes $\mu/\lambda$ constant, because
$Y/H$ is constant. Therefore $\dot\mu/\mu=\dot\lambda/\lambda=-\theta g$.
Using \eqref{eq:eg-human-costate} and the skill law gives
\begin{equation}\label{eq:eg-human-growth}
g=\frac{\eta-\delta_H-\rho}{\theta},\qquad
v^*=1-\frac{g+\delta_H}{\eta}
=\frac{\rho-(1-\theta)g}{\eta}.
\end{equation}
The stated restrictions ensure $0<v^*<1$. The physical Euler condition and
resource constraint also require
\begin{equation}
\alpha Y/K=\rho+\delta+\theta g,\qquad
C/K=Y/K-\delta-g>0.
\end{equation}
Since $Y/K=(vH/K)^{1-\alpha}$, the balanced stock ratio is
\begin{equation}
(K/H)^*=v^*\left(\frac{\alpha}{\rho+\delta+\theta g}\right)^{1/(1-\alpha)}.
\end{equation}
These equations characterize a balanced path, not the complete transition
from arbitrary $K_0/H_0$. Consumption and time allocation can adjust on impact;
the two stocks cannot be chosen afresh to match this ratio.

\Needspace{10\baselineskip}
\begin{redblock}
\textbf{Where did diminishing returns go?}
Adding machines while holding skill fixed still lowers their marginal product.
But scaling $K$ and $H$ together at a fixed work fraction scales output by
the same factor. The learning technology also reproduces $H$ linearly.
Persistent growth uses those joint reproduction opportunities, rather than
assuming that machines alone have escaped diminishing returns. If learning
became progressively harder as $H$ rose, this conclusion would need to be
rederived.
\end{redblock}

More schooling time at a given $H$ raises immediate skill accumulation while
reducing current goods production. A one-time higher $H_0$ is different from
a permanent improvement in training productivity $\eta$. School attendance,
actual learning, skill depreciation, and the use of skills at work are not
interchangeable measures. The model specifies one way those margins interact;
it does not imply that every extra school year permanently raises a country's
growth rate. Shared ideas introduce another issue: unlike a person's time,
the same design can be used in many places at once.

\subsection{Ideas, innovation, and incentives}
\textit{This section is being developed from the original outline.}
Separate nonrival ideas, private innovation rewards, and the social return.

\subsection{Allocation and aggregate productivity}
\textit{This section is being developed from the original outline.}
Connect the allocation of a given capital stock to productivity and growth.
